\section{Orçamentos rígidos e erro uniforme}
\label{sec:hard-budgets}

Além do funcional penalizado, um orçamento explícito de folhas pode ser
incorporado ao estado sem alterar o princípio de otimalidade.

\subsection{Número exato de folhas}

Para $1\le n\le2^d$, defina $C_{d,n}(J;u,v)$ como o menor valor de
\eqref{eq:additive-objective} entre representações com exatamente $n$ folhas.
A recorrência é
\begin{align}
 C_{0,1}(J;u,v)&=\Phi_H(J;u,v),\label{eq:budget-base}\\
 C_{0,n}(J;u,v)&=+\infty,\qquad n>1,\notag\\
 C_{d,1}(J;u,v)&=\Phi_H(J;u,v),\qquad d\ge1,\label{eq:budget-one}\\
 C_{d,n}(J;u,v)&=\lambda_{\rm bits}+\lambda_{\rm split}\notag\\
 &\quad+\min_{\substack{n_L+n_R=n\\n_L,n_R\ge1}}
 \min_{w\in\Values_{D,H}(K)}
 \{C_{d-1,n_L}(J_L;u,w)+C_{d-1,n_R}(J_R;w,v)\},
 \quad n\ge2.
 \label{eq:budget-recurrence}
\end{align}

\begin{theorem}[correção sob orçamento de folhas]
\label{thm:leaf-budget}
As recorrências \eqref{eq:budget-base}--\eqref{eq:budget-recurrence}
calculam exatamente o menor custo entre representações de profundidade no
máximo $d$, altura no máximo $H$, fronteiras $(u,v)$ e exatamente $n$
folhas. Em particular, o ótimo com no máximo $N$ folhas é
\begin{equation}
 \min_{\substack{1\le n\le\min\{N,2^D\}\\
                  u,v\in\Values_{D,H}(K)}}
 C_{D,n}(K;u,v).
 \label{eq:at-most-N}
\end{equation}
\end{theorem}

\begin{proof}
A indução em $d$ é idêntica à do
\cref{thm:variable-correctness}, com uma coordenada adicional. Uma
representação com uma folha usa o primeiro ramo. Se possui $n\ge2$ folhas,
sua raiz é interna e as quantidades positivas de folhas dos filhos satisfazem
$n_L+n_R=n$. Reciprocamente, toda escolha listada em
\eqref{eq:budget-recurrence} cola duas representações com exatamente $n$
folhas. A união sobre $1\le n\le N$ prova \eqref{eq:at-most-N}.
\end{proof}

\subsection{Recorrência de gargalo para a norma uniforme}

Seja agora $f\in C(K)$ e suponha disponível, para cada par admissível, o erro
local exato
\begin{equation}
 e_f(J;u,v)=\|f-\ell_{J;u,v}\|_{\infty,J}.
 \label{eq:local-uniform-error}
\end{equation}
Quando o par é inadmissível, ponha $e_f=+\infty$. Não se supõe que esse oráculo
exista para toda apresentação de uma função contínua; para polinômios e
famílias elementares ele pode ser produzido por cálculo simbólico ou por
aritmética intervalar certificada.

Defina $B_{d,n}(J;u,v)$ substituindo soma de perdas pelo máximo entre folhas:
\begin{align}
 B_{0,1}(J;u,v)&=e_f(J;u,v),\label{eq:bottleneck-base}\\
 B_{0,n}(J;u,v)&=+\infty,\qquad n>1,\notag\\
 B_{d,1}(J;u,v)&=e_f(J;u,v),\label{eq:bottleneck-one}\\
 B_{d,n}(J;u,v)
 &=\min_{\substack{n_L+n_R=n\\n_L,n_R\ge1}}
 \min_{w\in\Values_{D,H}(K)}
 \max\{B_{d-1,n_L}(J_L;u,w),
       B_{d-1,n_R}(J_R;w,v)\}.
 \label{eq:bottleneck-recurrence}
\end{align}

Seja
\begin{equation}
 E_{N,H,D}^{\rm tree}(f)_\infty
 =\min\{\|f-\Eval(R)\|_{\infty,K}:
 R\in\Representations_{D,H}^{\Library}(K),\
 |\Lambda(R)|\le N\}.
 \label{eq:truncated-tree-error}
\end{equation}

\begin{theorem}[cálculo exato do melhor erro truncado]
\label{thm:bottleneck}
Sob o oráculo \eqref{eq:local-uniform-error},
\begin{equation}
 E_{N,H,D}^{\rm tree}(f)_\infty
 =\min_{\substack{1\le n\le\min\{N,2^D\}\\
                  u,v\in\Values_{D,H}(K)}}
 B_{D,n}(K;u,v).
 \label{eq:bottleneck-optimum}
\end{equation}
Se
\begin{equation}
 D\ge \max\{0,H-k_0-1\},
 \label{eq:depth-inactive}
\end{equation}
a restrição de profundidade já é inativa: nenhum átomo de $\Library_H$
apoiado em descendente de $K$ pode ocorrer abaixo desse nível.
\end{theorem}

\begin{proof}
Em uma montagem contínua, o erro uniforme global é o máximo dos erros nas
folhas. A mesma partição em raiz-folha ou raiz-interna usada no
\cref{thm:leaf-budget} produz
\eqref{eq:bottleneck-recurrence}, agora com o operador máximo na colagem.
A indução em $d$ prova que $B_{d,n}$ é o melhor gargalo com o estado e a
quantidade de folhas fixados. Minimizar essas coordenadas dá
\eqref{eq:bottleneck-optimum}.

Uma folha na profundidade relativa $r$ tem escala $k_0+r$. Pela altura
\eqref{eq:inherited-height}, sua admissibilidade exige
$k_0+r+1\le H$. Portanto $r\le H-k_0-1$, o que prova a última afirmação.
\end{proof}

\begin{remark}[interação entre profundidade e translação]
\label{rem:depth-vs-translation}
A condição \eqref{eq:depth-inactive} satura o ganho obtido pelo aumento de $D$
além de $H-k_0-1$. A admissibilidade de cada célula também depende de sua
posição.
Para $m_0\ge0$, a \cref{prop:translation-obstruction} mostra que o refinamento
uniforme completo exige também
$r\le\log_2\bigl(H/(m_0+1)\bigr)$. Quando esse segundo limiar é o menor, o
catálogo se torna espacialmente assimétrico antes de a escala se esgotar. Ao ler
\eqref{eq:bottleneck-optimum} para um dado $(N,H,D)$, portanto, o orçamento de
folhas $N$ compete com uma família de árvores que já está truncada pela
translação, não apenas pela profundidade.
\end{remark}

\begin{corollary}[convergência para funções contínuas]
\label{cor:uniform-convergence}
Para todo $f\in C(K)$,
\begin{equation}
 \inf_{N,H,D}E_{N,H,D}^{\rm tree}(f)_\infty=0.
 \label{eq:uniform-convergence}
\end{equation}
Mais precisamente, ao longo de qualquer sequência de triplas $(N,H,D)$ que
seja componente a componente não decrescente e cofinal, os erros convergem a
zero.
\end{corollary}

\begin{proof}
Os catálogos são aninhados em cada orçamento, logo os erros não crescem. O
teorema de densidade da \cref{sec:assemblies} afirma que, dado $\varepsilon>0$,
existe uma montagem contínua finita de átomos com erro menor que
$\varepsilon$. Essa representação possui números finitos de folhas,
profundidade e altura; portanto pertence a algum catálogo da união. Assim, o
ínfimo em \eqref{eq:uniform-convergence} é zero, e a monotonicidade prova a
convergência cofinal.
\end{proof}
